%! Rates of Change in Polar Functions — Unit 7, Lesson 7.8 — Exit Ticket
\documentclass[10pt]{article}
\usepackage{precalculus-article}
\usepackage{precalculus-boxes}
\newcommand{\TallMath}[1]{$\displaystyle #1\rule[-1.4em]{0pt}{3.2em}$}
\newcommand{\CourseName}{Precalculus}
\newcommand{\MeetingLength}{55 minutes}
\newcommand{\UnitNumberName}{Unit 7: Analytic Trigonometry and Polar Coordinates \quad}
\newcommand{\LessonNumberName}{Lesson 7.8: Rates of Change in Polar Functions}

\begin{document}

\pageheader{Unit 7, Lesson 7.8}{Exit Ticket}
\namedateperiod

\vspace{0.12in}

\begin{notesbox}{Exit Ticket --- Answer independently}

\textbf{1.}\quad The table below gives partial values of $r = 1 + \sin\theta$.

\vspace{6pt}
\renewcommand{\arraystretch}{1.5}
\begin{tabular}{|c|c|c|c|c|}
\hline
$\theta$ & $0$ & $\pi/2$ & $\pi$ & $3\pi/2$ \\
\hline
$r = 1 + \sin\theta$ & $1$ & $2$ & $1$ & $0$ \\
\hline
\end{tabular}

\vspace{8pt}
(a)\quad On the interval $[0, \pi/2]$, is $r$ \textbf{positive and increasing}
or \textbf{positive and decreasing}?

\vspace{4pt}
\underline{\hspace{7cm}}

\vspace{8pt}
(b)\quad On the interval $[\pi/2, \pi]$, is $r$ \textbf{positive and increasing}
or \textbf{positive and decreasing}?

\vspace{4pt}
\underline{\hspace{7cm}}

\vspace{0.16in}
\textbf{2.}\quad Compute the average rate of change of $r = 1 + \sin\theta$
from $\theta = \pi/2$ to $\theta = \pi$.

\vspace{4pt}
$$\frac{r(\pi) - r(\pi/2)}{\pi - \pi/2} = \frac{\rule{1.5cm}{0.4pt} - \rule{1.5cm}{0.4pt}}{\rule{2cm}{0.4pt}} = \underline{\hspace{3.5cm}}$$

\vspace{0.16in}
\textbf{3.}\quad In one sentence, explain what it means for a polar function to have
a \textbf{negative} average rate of change of $r$ over an interval.

\writelines{2}

\end{notesbox}

\end{document}
